\documentclass{amsart}
% For dealing with references we use the comment environment
\usepackage{verbatim}
\newenvironment{reference}{\comment}{\endcomment}
%\newenvironment{reference}{}{}
\newenvironment{slogan}{\comment}{\endcomment}
\newenvironment{history}{\comment}{\endcomment}

% For commutative diagrams we use Xy-pic
\usepackage[all]{xy}

% We use 2cell for 2-commutative diagrams.
\xyoption{2cell}
\UseAllTwocells

% We use multicol for the list of chapters between chapters
\usepackage{multicol}

% This is generally recommended for better output
\usepackage{lmodern}
\usepackage[T1]{fontenc}

% For cross-file-references

% Package for hypertext links:
\usepackage{hyperref}

% For any local file, say "hello.tex" you want to link to please

% Theorem environments.
%
\theoremstyle{plain}
\newtheorem{theorem}[subsection]{Theorem}
\newtheorem{proposition}[subsection]{Proposition}
\newtheorem{lemma}[subsection]{Lemma}

\theoremstyle{definition}
\newtheorem{definition}[subsection]{Definition}
\newtheorem{example}[subsection]{Example}
\newtheorem{exercise}[subsection]{Exercise}
\newtheorem{situation}[subsection]{Situation}

\theoremstyle{remark}
\newtheorem{remark}[subsection]{Remark}
\newtheorem{remarks}[subsection]{Remarks}

\numberwithin{equation}{subsection}

% Macros
%
\def\lim{\mathop{\mathrm{lim}}\nolimits}
\def\colim{\mathop{\mathrm{colim}}\nolimits}
\def\Spec{\mathop{\mathrm{Spec}}}
\def\Hom{\mathop{\mathrm{Hom}}\nolimits}
\def\Ext{\mathop{\mathrm{Ext}}\nolimits}
\def\SheafHom{\mathop{\mathcal{H}\!\mathit{om}}\nolimits}
\def\SheafExt{\mathop{\mathcal{E}\!\mathit{xt}}\nolimits}
\def\Sch{\mathit{Sch}}
\def\Mor{\mathop{\mathrm{Mor}}\nolimits}
\def\Ob{\mathop{\mathrm{Ob}}\nolimits}
\def\Sh{\mathop{\mathit{Sh}}\nolimits}
\def\NL{\mathop{N\!L}\nolimits}
\def\CH{\mathop{\mathrm{CH}}\nolimits}
\def\proetale{{pro\text{-}\acute{e}tale}}
\def\etale{{\acute{e}tale}}
\def\QCoh{\mathit{QCoh}}
\def\Ker{\mathop{\mathrm{Ker}}}
\def\Im{\mathop{\mathrm{Im}}}
\def\Coker{\mathop{\mathrm{Coker}}}
\def\Coim{\mathop{\mathrm{Coim}}}

% Boxtimes
%
\DeclareMathSymbol{\boxtimes}{\mathbin}{AMSa}{"02}

%
% Macros for moduli stacks/spaces
%
\def\QCohstack{\mathcal{QC}\!\mathit{oh}}
\def\Cohstack{\mathcal{C}\!\mathit{oh}}
\def\Spacesstack{\mathcal{S}\!\mathit{paces}}
\def\Quotfunctor{\mathrm{Quot}}
\def\Hilbfunctor{\mathrm{Hilb}}
\def\Curvesstack{\mathcal{C}\!\mathit{urves}}
\def\Polarizedstack{\mathcal{P}\!\mathit{olarized}}
\def\Complexesstack{\mathcal{C}\!\mathit{omplexes}}
% \Pic is the operator that assigns to X its picard group, usage \Pic(X)
% \Picardstack_{X/B} denotes the Picard stack of X over B
% \Picardfunctor_{X/B} denotes the Picard functor of X over B
\def\Pic{\mathop{\mathrm{Pic}}\nolimits}
\def\Picardstack{\mathcal{P}\!\mathit{ic}}
\def\Picardfunctor{\mathrm{Pic}}
\def\Deformationcategory{\mathcal{D}\!\mathit{ef}}

\setlength{\emergencystretch}{3em}
\begin{document}
\title[Stacks Project, Tag 0BP2]{410 --- Beauville-Laszlo gluing equivalence for modules}
\maketitle
\noindent Copyright (C) 2005--2025 Johan de Jong. Stacks Project authors. Source: \href{https://stacks.math.columbia.edu/tag/0BP2}{Tag 0BP2}.

\noindent Editorial difficulty note (not author text). Difficulty H3: The proof constructs the glued module as the kernel of a surjective difference map. It separates power torsion from the torsion-free quotient, controls Tor under base change and uses a comparison diagram to recover both local pieces, proving full faithfulness and essential surjectivity..

\noindent Editorial provenance note (not author text). History: excerpt from revision \texttt{a0444\allowbreak 6e57e\allowbreak c1fbc\allowbreak 252a8\allowbreak 71afc\allowbreak ec775\allowbreak 2fb28\allowbreak 07b14}, more-algebra.tex, lines 25578--25673. Reference presentation was adapted; original-author context and core support blocks were retained or added. The mathematical wording is unchanged. Licensed under GNU FDL 1.2 or later; no invariant sections or cover texts. See ../licenses/COPYING. Context blocks reproduce author definitions and supporting results; they are not additional counted problems.

{\bf Glueing pairs.} Let $R \to R'$ be a ring map that induces isomorphisms
$R/f^nR \to R'/f^nR'$ for $n > 0$. Consider the sequence
\begin{equation}
\label{equation-BL-cech-re}
0 \to R \to R' \oplus R_f \to R'_f \to 0,
\end{equation}
in which the map on the right is the difference between the two canonical
homomorphisms. If this sequence is exact, then
we say that $(R \to R', f)$ is a \emph{glueing pair}.
We will say that $(R, f)$ is a \emph{glueing pair} if
$(R \to R^\wedge, f)$ is a glueing pair; this makes sense by
Lemma \href{https://stacks.math.columbia.edu/tag/0BNJ}{lemma same quotients (Tag 0BNJ)}. Thus $(R, f)$
is a glueing pair if and only if the sequence
\begin{equation}
\label{equation-BL-cech}
0 \to R \to R^\wedge \oplus R_f \to (R^\wedge)_f \to 0,
\end{equation}
is exact.



{\bf Glueable modules.}
Let $R \to R'$ be a ring map which induces isomorphisms
$R/f^nR \to R'/f^nR'$ for $n > 0$. For any $R$-module $M$, we
may tensor (\ref{equation-BL-cech-re}) with $M$ to obtain a sequence
\begin{equation}
\label{equation-BL-cech-mod-re}
0 \to M \to (M \otimes_R R') \oplus (M \otimes_R R_f) \to
M \otimes_R R'_f \to 0
\end{equation}
Observe that $M \otimes_R R_f = M_f$ and that
$M \otimes_R R'_f = (M \otimes_R R')_f$.
If this sequence is exact, we say that $M$ is
\emph{glueable for $(R \to R', f)$}.
If $R$ is a ring and $f \in R$, then we say an $R$-module
is \emph{glueable} if $M$ is glueable for $(R \to R^\wedge, f)$.
Thus $M$ is glueable if and only if the sequence
\begin{equation}
\label{equation-BL-cech-mod}
0 \to M \to (M \otimes_R R^\wedge) \oplus (M \otimes_R R_f) \to
M \otimes_R (R^\wedge)_f \to 0
\end{equation}
is exact.



\begin{remark}
\label{remark-glueing-data}
In this remark we define a category of glueing data.
Let $R \to S$ be a ring map.
Let $f_1, \ldots, f_t \in R$ and $I = (f_1, \ldots, f_t)$.
Consider the category $\text{Glue}(R \to S, f_1, \ldots, f_t)$
as the category whose
\begin{enumerate}
\item objects are systems $(M', M_i, \alpha_i, \alpha_{ij})$, where
$M'$ is an $S$-module, $M_i$ is an $R_{f_i}$-module,
$\alpha_i : (M')_{f_i} \to M_i \otimes_R S$ is an isomorphism, and
$\alpha_{ij} : (M_i)_{f_j} \to (M_j)_{f_i}$ are isomorphisms
such that
\begin{enumerate}
\item $\alpha_{ij} \circ \alpha_i = \alpha_j$ as maps
$(M')_{f_if_j} \to (M_j)_{f_i}$, and
\item $\alpha_{jk} \circ \alpha_{ij} = \alpha_{ik}$ as maps
$(M_i)_{f_jf_k} \to (M_k)_{f_if_j}$ (cocycle condition).
\end{enumerate}
\item morphisms
$(M', M_i, \alpha_i, \alpha_{ij}) \to (N', N_i, \beta_i, \beta_{ij})$
are given by maps $\varphi' : M' \to N'$ and $\varphi_i : M_i \to N_i$
compatible with the given maps $\alpha_i, \beta_i, \alpha_{ij}, \beta_{ij}$.
\end{enumerate}
There is a canonical functor
$$
\text{Can} : \text{Mod}_R
\longrightarrow
\text{Glue}(R \to S, f_1, \ldots, f_t),
\quad
M \longmapsto (M \otimes_R S, M_{f_i}, \text{can}_i, \text{can}_{ij})
$$
where $\text{can}_i : (M \otimes_R S)_{f_i} \to M_{f_i} \otimes_R S$
and $\text{can}_{ij} : (M_{f_i})_{f_j} \to (M_{f_j})_{f_i}$
are the canonical isomorphisms. For any object
$\mathbf{M} = (M', M_i, \alpha_i, \alpha_{ij})$ of the category
$\text{Glue}(R \to S, f_1, \ldots, f_t)$ we define
$$
H^0(\mathbf{M}) =
\{(m', m_i) \mid \alpha_i(m') = m_i \otimes 1, \alpha_{ij}(m_i) = m_j\}
$$
in other words defined by the exact sequence
$$
0 \to H^0(\mathbf{M}) \to
M' \times \prod M_i \to
\prod M'_{f_i}
\times
\prod (M_i)_{f_j}
$$
similar to (\href{https://stacks.math.columbia.edu/tag/05EJ}{equation glueing complex (Tag 05EJ)}).
We think of $H^0(\mathbf{M})$ as an $R$-module. Thus we also get a functor
$$
H^0 :
\text{Glue}(R \to S, f_1, \ldots, f_t)
\longrightarrow
\text{Mod}_R
$$
Our next goal is to show that the functors
$\text{Can}$ and $H^0$ are sometimes quasi-inverse to each other.
\end{remark}


\begin{lemma}
\label{lemma-BL3}
\begin{reference}
Slight generalization of \href{https://stacks.math.columbia.edu/bibliography}{Bibliography: Beauville-Laszlo (Lemme 3(a))}
\end{reference}
Let $(R \to R', f)$ be a glueing pair. For every $R$-module $M$, we have
$\text{Tor}^R_1(R', \Coker(M \to M_f)) = 0$.
\end{lemma}

\begin{proof}
Set $\overline{M} = M/M[f^\infty]$. Then
$\Coker(M \to M_f) \cong \Coker(\overline{M} \to \overline{M}_f)$
hence we may and do assume that $f$ is a nonzerodivisor on $M$.
In this case $M \subset M_f$ and $M_f/M = \colim M/f^nM$ where the
transition maps are given by multiplication by $f$. Since
formation of Tor groups commutes with colimits
(Algebra, Lemma \href{https://stacks.math.columbia.edu/tag/0BNF}{lemma tor commutes filtered colimits (Tag 0BNF)})
it suffices to show that $\text{Tor}^R_1(R', M/f^n M) = 0$.

\medskip\noindent
We first treat the case $M = R/R[f^\infty]$. By
Lemma \href{https://stacks.math.columbia.edu/tag/0BNR}{lemma same f torsion (Tag 0BNR)}
we have $M \otimes_R R' = R'/R'[f^\infty]$.
From the short exact sequence $0 \to M \to M \to M/f^nM \to 0$
we obtain the exact sequence
$$
\xymatrix{
\text{Tor}_1^R(R', R/R[f^\infty]) \ar[r] &
\text{Tor}_1^R(R', M/f^n M) \ar[r] &
R'/R'[f^\infty] \ar[dll]_{f^n} \\
R'/R'[f^\infty] \ar[r] &
(R'/R'[f^\infty])/(f^n
(R'/R'[f^\infty])) \ar[r] & 0
}
$$
by Algebra, Lemma \href{https://stacks.math.columbia.edu/tag/00M0}{lemma long exact sequence tor (Tag 00M0)}.
Here the diagonal arrow is injective. Since the first group
$\text{Tor}_1^R(R', R/R[f^\infty])$ is zero by
Lemma \href{https://stacks.math.columbia.edu/tag/0BP0}{lemma first tor total (Tag 0BP0)}, we deduce that
$\text{Tor}_1^R(R', M/f^nM) = 0$ as desired.

\medskip\noindent
To treat the general case, choose a surjection $F \to M$ with $F$ a free
$R/R[f^\infty]$-module, and form an exact sequence
$$
0 \to N \to F/f^n F \to M/f^n M \to 0.
$$
By Lemma \href{https://stacks.math.columbia.edu/tag/0BNK}{lemma torsion completion (Tag 0BNK)}
this sequence remains unchanged, and hence
exact, upon tensoring with $R'$.
Since $\text{Tor}^R_1(R', F/f^n F) = 0$ by the
previous paragraph, we deduce  that
$\text{Tor}^R_1(R', M/f^n M) = 0$ as desired.
\end{proof}

\begin{theorem}
\label{theorem-BL-glueing}
\begin{reference}
Slight generalization of the main theorem of \href{https://stacks.math.columbia.edu/bibliography}{Bibliography: Beauville-Laszlo}.
\end{reference}
Let $(R \to R',f)$ be a glueing pair. The functor
$\text{Can} : \text{Mod}_R \longrightarrow \text{Glue}(R \to R', f)$
determines an equivalence of the category of $R$-modules glueable
for $(R \to R', f)$ and the category $\text{Glue}(R \to R', f)$
of glueing data.
\end{theorem}

\begin{proof}
Let $(M', M_1, \alpha_1)$ be a glueing datum. We will show that
$M = H^0((M', M_1, \alpha_1))$ is a glueable module for $(R \to R', f)$
and that $(M', M_1, \alpha_1) \cong \text{Can}(M)$.

\medskip\noindent
We first check that the map $\text{d} : M' \oplus M_1 \to (M')_f$ used in the
definition of the functor $H^0$ is surjective. Observe that
$(x, y) \in M' \oplus M_1$ maps to
$\text{d}(x, y) = x/1 - \alpha_1^{-1}(y \otimes 1)$
in $(M')_f$. If $z \in (M')_f$, then we can write
$\alpha_1(z) = \sum y_i \otimes g_i$ with $g_i \in R'$
and $y_i \in M_1$. Write $\alpha_1^{-1}(y_i \otimes 1) = y_i'/f^n$
for some $y'_i \in M'$ and $n \geq 0$ (we can pick the same $n$
for all $i$). Write $g_i = a_i + f^n b_i$ with $a_i \in R$ and
$b_i \in R'$. Then with $y = \sum a_i y_i \in M_1$ and
$x = \sum b_i y'_i \in M'$ we have $\text{d}(x, -y) = z$
as desired.

\medskip\noindent
Since $M = H^0((M', M_1, \alpha_1)) = \Ker(\text{d})$ we obtain
an exact sequence of $R$-modules
\begin{equation}
\label{equation-define-M}
0 \to M \to M' \oplus M_1 \to (M')_f \to 0.
\end{equation}
We will prove that the maps $M \to M'$ and $M \to M_1$ induce isomorphisms
$M \otimes_R R' \to M'$ and $M \otimes_R R_f \to M_1$.
This will imply that $M$ is glueable for $(R \to R', f)$ and
$\text{Can}(M) \cong (M', M_1, \alpha_1)$ as desired.

\medskip\noindent
Since $f$ is a nonzerodivisor on $M_1$, we have
$M[f^\infty] \cong M'[f^\infty]$. This yields an exact sequence
\begin{equation}
\label{equation-exact-mod-torsion}
0 \to M/M[f^\infty] \to M_1 \to (M')_f/M' \to 0.
\end{equation}
Since $R \to R_f$ is flat, we may tensor this exact sequence with $R_f$
to deduce that $M \otimes_R R_f = (M/M[f^\infty]) \otimes_R R_f \to M_1$
is an isomorphism.

\medskip\noindent
By Lemma \ref{lemma-BL3} we have
$\text{Tor}_1^R(R', \Coker(M' \to (M')_f)) = 0$.
The sequence (\ref{equation-exact-mod-torsion})
thus remains exact upon tensoring
over $R$ with $R'$. Using $\alpha_1$ and
Lemma \href{https://stacks.math.columbia.edu/tag/0BNK}{lemma torsion completion (Tag 0BNK)}
the resulting exact sequence can be written as
\begin{equation}
\label{equation-mod-torsion-sequence}
0 \to (M/M[f^\infty]) \otimes_R R' \to
(M')_f \to (M')_f/M' \to 0
\end{equation}
This yields an isomorphism
$(M/M[f^\infty]) \otimes_R R' \cong M'/M'[f^\infty]$.
This implies that in the diagram
$$
\xymatrix{
& M[f^\infty] \otimes_R R' \ar[r] \ar[d] &
M \otimes_R R'  \ar[r] \ar[d] &
(M/M[f^\infty]) \otimes_R R' \ar[r] \ar[d] & 0 \\
0 \ar[r] &
M'[f^\infty] \ar[r] &
M' \ar[r] &
M'/M'[f^\infty] \ar[r] & 0,
}
$$
the third vertical arrow is an isomorphism. Since the rows are exact and the
first vertical arrow is an isomorphism by
Lemma \href{https://stacks.math.columbia.edu/tag/0BNK}{lemma torsion completion (Tag 0BNK)} and $M[f^\infty] = M'[f^\infty]$,
the five lemma implies that $M \otimes_R R' \to M'$ is an isomorphism.

\medskip\noindent
The above shows that $\text{Can}$ is essentially surjective and that
the functor $H^0$ maps into the category of glueable modules.
Due to the exactness of (\ref{equation-BL-cech-mod-re}) for glueable modules
we have $H^0 \circ \text{Can} = \text{id}$ on the
category of glueable modules. This implies $\text{Can}$ is fully
faithful by Lemma \href{https://stacks.math.columbia.edu/tag/0H77}{lemma H0 Can adjoint (Tag 0H77)} combined with
Categories, Lemma \href{https://stacks.math.columbia.edu/tag/07RB}{lemma adjoint fully faithful (Tag 07RB)}.
This finishes the proof.
\end{proof}
\end{document}
